\subsection{Instrumentación electrónica}\label{sec:montaje_instrumentacion}

La instrumentación del prototipo se organiza alrededor del PLC Siemens S7-1200 CPU 1214C AC/DC/RLY, que recibe las señales de los sensores del tanque de mezcla y acciona los actuadores a través de sus salidas de relé, directamente o mediante relés de estado sólido (SSR) y relés auxiliares. La Tabla~\ref{tab:instrumentacion} resume los elementos instalados, con la nomenclatura del P\&ID.

\begin{table}[H]
\centering
\footnotesize
\caption{Sensores y actuadores instalados en el prototipo.}
\label{tab:instrumentacion}
\begin{tabularx}{\textwidth}{L{1.4cm}L{3.6cm}YL{3.6cm}}
\toprule
\textbf{Tag} & \textbf{Elemento} & \textbf{Descripción} & \textbf{Conexión con el PLC}\\
\midrule
TT-101 & Transmisor de temperatura & Sensor RTD PT100 de 3 hilos con transmisor SITRANS de 4--20\,mA, instalado en el tanque de mezcla & Entrada analógica, mediante una resistencia de 500\,$\Omega$; escalamiento con \texttt{NORM\_X} y \texttt{SCALE\_X}\\
LSL-101 / LSH-101 & Interruptores de nivel bajo y alto & Flotador doble instalado en el tanque de mezcla & Entradas digitales\\
E-101 & Resistencia de inmersión & Resistencia de 1500\,W de placa, que entrega aproximadamente 2000\,W a la tensión del laboratorio (126\,V) & Salida de relé que activa un SSR de 60\,A\\
P-102 & Bomba de impulsión & Bomba periférica de 110\,V\,AC & Salida de relé que activa un SSR de 40\,A\\
P-103 & Bomba de agitación & Bomba sumergible de 22\,W y 110\,V\,AC, ubicada dentro del tanque de mezcla & Salida de relé que activa un relé auxiliar con bobina de 24\,V\\
P-101 & Bomba dosificadora & Bomba peristáltica de 12\,V\,DC, alimentada mediante un convertidor reductor LM2596 & Salida de relé (encendido y apagado temporizado)\\
SV-102 & Electroválvula de llenado & Electroválvula de 24\,V\,DC en la entrada de agua de red & Salida de relé\\
SV-103 & Electroválvula de descarga & Electroválvula de 24\,V\,DC entre el tanque de mezcla y la bomba de impulsión & Salida de relé\\
--- & Elementos de mando & Paro de emergencia tipo hongo, pulsadores, selector y pilotos luminosos & Entradas y salidas digitales\\
\bottomrule
\end{tabularx}
\end{table}

\paragraph{Alimentación.} El PLC se alimenta a 110\,V\,AC desde la red. Los 24\,V\,DC para los sensores, el transmisor de temperatura y los circuitos de mando provienen de una fuente WDR-120-24 (24\,V, 120\,W) montada en el tablero, y no de la salida de sensores del PLC. En el caso de las electroválvulas, el contacto de la salida de relé del PLC conmuta el positivo (+24\,V) de esta fuente hacia la bobina de cada electroválvula, cuyo otro terminal se conecta al negativo (0\,V) de la misma fuente.

\paragraph{Protecciones.} La resistencia de inmersión cuenta con un interruptor automático de 15\,A de un polo, y la bomba de impulsión y la bomba de agitación comparten un interruptor automático independiente de 7\,A. El sistema dispone además de un paro de emergencia tipo hongo.

\paragraph{Montaje y conexionado.} Los disyuntores, la fuente de alimentación y el transmisor de temperatura se ubicaron en un tablero elaborado por el equipo. Las conexiones de sensores y actuadores se realizaron con tornillería métrica y se protegieron con termoencogible, de modo que quedaran fijas, seguras y aisladas de la humedad.

\campo{Pendiente: diagrama unifilar de potencia, diagrama de conexiones del PLC (mapa de entradas y salidas con sus direcciones) y fotografías del tablero.}
